\documentclass[10pt]{article}
\usepackage[margin=0.85in]{geometry}
\usepackage{amsmath,amssymb,amsthm}
\usepackage[hidelinks]{hyperref}
\newtheorem{theorem}{Theorem}
\title{Minimal left ideals are not matrix factors\\Conjecture 00000007627}
\author{earthking11}
\date{October 6, 2026}
\begin{document}
\maketitle
\section*{Statement and interpretation}
The source defines ideals as left-module structures and asserts that their
minimal-ideal count equals the matrix-factor count, with a center-dimension
counting law. It gives no ring or field restriction. We use the usual
meaning of minimal nonzero left ideals and matrix factors in a finite
direct-product decomposition. We count distinct ideals, not their module
isomorphism classes. The latter is a different assertion.

\begin{theorem}
Let $K=\mathbb F_2$ and $R=M_2(K)$, considered as its own left regular
module. The actual finite set of minimal left ideals has cardinality at
least $2$. Nevertheless $R$ is a single simple matrix factor, and its
center has dimension $1$ over $K$. Thus the two claimed counts disagree.
\end{theorem}
\begin{proof}
Consider the two column-supported ideals
\[
 L_0=\left\{\begin{pmatrix}a&0\\b&0\end{pmatrix}:a,b\in K\right\},
 \qquad
 L_1=\left\{\begin{pmatrix}0&a\\0&b\end{pmatrix}:a,b\in K\right\}.
\]
They are closed under addition and left multiplication by arbitrary
matrices in $R$. Each is nonzero, and $E_{00}$ belongs to $L_0$ but not
$L_1$, so they are distinct.

To prove minimality, take a nonzero matrix $A\in L_0$. At least one entry
in its first column is $1$. Let its row index be $j$. For any $B\in L_0$,
choose $C$ whose column $j$ is the first column of $B$, with its other
column zero. Then $CA=B$. Hence any nonzero left ideal inside $L_0$
contains all of $L_0$. This proves minimality. The same argument, using
the second column, proves that $L_1$ is minimal.

Since $R$ is finite, its set of submodules and the subset of minimal
left ideals are finite. The two distinct ideals give an injection from
a two-element set, so the actual minimal-left-ideal count is at least $2$.
On the other hand, $R$ is a simple matrix ring and has the genuine
one-factor decomposition
\[
 R\ \cong\ \prod_{i\in\{0\}} M_2(K),
 \qquad A\longmapsto (i\longmapsto A).
\]
The inverse evaluates at the sole index, and the actual factor-index
set has cardinality $1$. Finally, the center of $M_2(K)$ is the scalar
copy of $K$, so its dimension over $K$ is $1$. Neither count can equal
the minimal-left-ideal count.
\end{proof}

\section*{Formal verification and limits}
Lean uses the actual matrix ring, its left regular module
\texttt{Submodule R R}, the subtype of all minimal left ideals, and the
actual \texttt{Subalgebra.center}. It proves minimality against every
nonzero subideal, finite cardinality, the injection, the one-factor ring
equivalence, ring simplicity and the center dimension. The capstone
combines the decomposition witness with both actual counting mismatches.

The project builds on Lean 4.33.1 with pinned Mathlib; printed audits use
only propositional extensionality, classical choice and quotient soundness.
No unfinished proof, native evaluation, extra axiom or long search is used.
The source does not restrict this problem to Clifford algebras; no claim
that this example is a Clifford algebra is needed or made. We do not
assign a new meaning to its unspecified ``direct-sum index.''
\end{document}
